\section*{Bab \ref{ch:b2:hermitian} --- Bentuk Hermitian}

\begin{solution}{exo:b2:hermitian:1}
$\langle x, y\rangle = \conj{1}\cdot\iu + \conj{\iu}\cdot 1 = \iu -
\iu = 0$: jadi ortogonal. Lalu $\norm x = \norm y = \sqrt{1 + 1} = \sqrt2$.
Basis ortonormalnya: $\bigl(\frac{1}{\sqrt2}(1, \iu),\;
\frac{1}{\sqrt2}(\iu, 1)\bigr)$ --- yakni pasangan ternormalkan itu sendiri.
\end{solution}

\begin{solution}{exo:b2:hermitian:2}
Yang pertama: ia sama dengan \omterm{def:b2:linalg:transpose}{transpos} sekawannya (sebab diagonalnya real dan
$\conj{\iu} = -\iu$ tertukar): jadi \omterm{def:b2:hermitian:adjoint}{Hermitian} (sehingga normal); tetapi tak
\omterm{def:b2:hermitian:adjoint}{uniter} (sebab $A^\dagger A \neq I$: kolomnya tak satuan).

Yang kedua: $A^\dagger A = \frac12\begin{pmatrix} 1 & -\iu\\ -\iu &
1\end{pmatrix}\begin{pmatrix} 1 & \iu\\ \iu & 1\end{pmatrix} =
\frac12\begin{pmatrix} 2 & 0\\ 0 & 2\end{pmatrix} = I$: jadi \omterm{def:b2:hermitian:adjoint}{uniter}
(sehingga normal); tetapi tak \omterm{def:b2:hermitian:adjoint}{Hermitian}.

Yang ketiga: $A^\dagger A = E_{22} \neq E_{11} = AA^\dagger$: jadi tak normal
(sehingga bukan \omterm{def:b2:hermitian:adjoint}{Hermitian} maupun \omterm{def:b2:hermitian:adjoint}{uniter}) --- yakni contoh tandingan
nilpoten yang baku.
\end{solution}

\begin{solution}{exo:b2:hermitian:3}
Ketunggalannya: $A = H + \iu K$ dengan $H^\dagger = H$ dan $K^\dagger = K$
memaksa $A^\dagger = H - \iu K$, jadi $H = \frac{A + A^\dagger}{2}$ dan
$K = \frac{A - A^\dagger}{2\iu}$; sedangkan rumus itu bersifat \omterm{def:b2:hermitian:adjoint}{Hermitian}
(periksa: $\bigl(\frac{A - A^\dagger}{2\iu}\bigr)^\dagger = \frac{
A^\dagger - A}{-2\iu} = K$) dan membangun ulang $A$: itulah keberadaannya.

Kenormalannya: $A^\dagger A - AA^\dagger = (H - \iu K)(H + \iu K) - (H
+ \iu K)(H - \iu K) = 2\iu(HK - KH)$: jadi ia lenyap bila dan hanya bila $HK = KH$.
\end{solution}

\begin{solution}{exo:b2:hermitian:4}
$\chi_A = (X-2)^2 - 1$: jadi \omterm{def:b2:reduction:eigen}{nilai eigennya} $3$ dan $1$. \omterm{def:b2:reduction:eigen}{Vektor eigennya}, lewat
perhitungan langsung:
\[
A\begin{pmatrix}1\\ -\iu\end{pmatrix}
= \begin{pmatrix} 2 + \iu(-\iu)\\ -\iu + 2(-\iu)\end{pmatrix}
= \begin{pmatrix} 3\\ -3\iu \end{pmatrix}
= 3\begin{pmatrix}1\\ -\iu\end{pmatrix},
\qquad
A\begin{pmatrix}1\\ \iu\end{pmatrix}
= \begin{pmatrix} 2 + \iu\cdot\iu\\ -\iu + 2\iu\end{pmatrix}
= \begin{pmatrix}1\\ \iu\end{pmatrix}.
\]
Basis eigen ortonormalnya: $u_1 = \frac{1}{\sqrt2}(1, -\iu)$
(bernilai eigen $3$) dan $u_2 = \frac{1}{\sqrt2}(1, \iu)$ (bernilai eigen
$1$); dengan keortogonalannya seperti pada \cref{exo:b2:hermitian:1}. Adapun pangkatnya,
lewat proyeksi spektralnya $A^k = 3^k u_1u_1^\dagger + 1^k\,
u_2u_2^\dagger$:
\[
A^k = U\begin{pmatrix} 3^k & 0\\ 0 & 1\end{pmatrix}U^\dagger
= \frac{3^k}{2}\begin{pmatrix} 1 & \iu\\ -\iu & 1\end{pmatrix}
+ \frac{1}{2}\begin{pmatrix} 1 & -\iu\\ \iu & 1\end{pmatrix}
= \frac12\begin{pmatrix}
3^k + 1 & (3^k - 1)\iu\\
-(3^k-1)\iu & 3^k + 1
\end{pmatrix}.
\]
(Periksa $k = 1$: ia memulihkan $A$.)
\end{solution}

\begin{solution}{exo:b2:hermitian:5}
\omterm{def:b2:metric:compact}{Kompak}: sebab tertutup (yakni prapeta $I$ oleh pemetaan \omterm{def:b2:metric:continuity}{kontinu} $U \mapsto
U^\dagger U$) dan terbatas (sebab kolomnya vektor satuan, jadi entrinya
bermodulus $\leq 1$) di dalam $\mathcal{M}_n(\C) \simeq \R^{2n^2}$.

\omterm{def:b2:reduction:eigen}{Nilai eigennya}: sebab $\operatorname{diag}(\lambda, 1, \dots, 1)$ bersifat \omterm{def:b2:hermitian:adjoint}{uniter}
untuk sembarang $\abs\lambda = 1$. \omterm{def:b2:linalg:det}{Determinannya}: $\abs{\det U}^2 = \det
U^\dagger \det U = \det(U^\dagger U) = 1$ (dengan memakai $\det A^\dagger =
\conj{\det A}$): jadi $\det U$ terletak pada lingkaran satuan.
\end{solution}

\begin{solution}{exo:b2:hermitian:6}
Menurut teorema spektral, bekerjalah dalam basis eigen ortonormal $H$:
sebab segalanya menyusut menjadi skalar $\lambda \in \R$ (yakni \omterm{def:b2:reduction:eigen}{nilai eigennya}).
Matriks $H + \iu I$ bernilai eigen $\lambda + \iu \neq 0$: jadi terbalikkan. Lalu $U$
bernilai eigen $\mu = \frac{\lambda - \iu}{\lambda + \iu}$, yang
bermodulus $1$ (sebab $\abs{\lambda - \iu} = \abs{\lambda + \iu}$ untuk
$\lambda$ yang real): jadi $U^\dagger U = I$ berlaku karena $U$ \omterm{def:b2:reduction:diag}{dapat
didiagonalkan} secara \omterm{def:b2:hermitian:adjoint}{uniter} dengan \omterm{def:b2:reduction:eigen}{nilai eigen} bermodulus satu (ia diagonal dalam
basis ortonormal yang dipilih). Dan $\mu = 1$ akan memaksa $-\iu = \iu$:
yang mustahil, jadi $1 \notin \operatorname{Sp} U$. (Transformasi Cayley
memetakan yang \omterm{def:b2:hermitian:adjoint}{Hermitian} ke yang \omterm{def:b2:hermitian:adjoint}{uniter} dikurangi satu titik --- yakni versi
matriks bagi pemetaan dari $\R$ ke lingkarannya.)
\end{solution}

\begin{solution}{exo:b2:hermitian:7}
Untuk pasangan eigen $Ax = \lambda x$ (dengan $x \neq 0$): $\lambda\norm x^2 =
\langle x, Ax\rangle > 0$, jadi $\lambda > 0$ (dan sudah real menurut
\cref{prop:b2:hermitian:eigenvalues}). Akar kuadratnya: dalam basis
spektral, $B = U\operatorname{diag}(\sqrt{\lambda_i})U^\dagger$ bersifat
\omterm{def:b2:hermitian:adjoint}{Hermitian}, definit positif, dan $B^2 = A$. \omterm{def:b2:linalg:det}{Determinannya}: yakni hasil kali
\omterm{def:b2:reduction:eigen}{nilai eigennya} yang positif.
\end{solution}

\begin{solution}{exo:b2:hermitian:8}
\begin{enumerate}
  \item $\norm{u(x)}^2 = \langle u(x), u(x)\rangle = \langle x,
        u^*u(x)\rangle = \langle x, uu^*(x)\rangle =
        \norm{u^*(x)}^2$. Menerapkan ini pada $u -
        \lambda\,\mathrm{id}$ yang normal (sebab \omterm{def:b2:hermitian:adjoint}{adjoinnya} $u^* -
        \conj\lambda$, dan kenormalannya diwarisi): $\norm{(u -
        \lambda)x} = \norm{(u^* - \conj\lambda)x}$, jadi kernelnya
        sepakat.
  \item Untuk \omterm{def:b2:reduction:eigen}{vektor eigen} $u(x) = \lambda x$ dan $u(y) = \mu y$
        (dengan $\lambda \neq \mu$): memakai (1), $u^*(x) = \conj\lambda x$;
        lalu
        \[
        \lambda\langle y, x\rangle = \langle y, u(x)\rangle
        = \langle u^*(y), x\rangle = \langle \conj\mu\, y, x\rangle
        = \mu \langle y, x\rangle ,
        \]
        jadi $\langle y, x\rangle = 0$. Adapun kestabilan $E_\lambda^\perp$:
        untuk $x \perp E_\lambda$ dan $z \in E_\lambda$, $\langle z,
        u(x)\rangle = \langle u^*(z), x\rangle =
        \langle\conj\lambda z, x\rangle = 0$.
  \item Secara induktif atas dimensinya: atas $\C$, $u$ punya \omterm{def:b2:reduction:eigen}{vektor eigen}
        $e_1$ (normalkan ia); lalu komplemen ortogonalnya stabil
        terhadap $u$ (menurut (2)) \emph{sekaligus} terhadap $u^*$ (lewat hujah
        yang sama dengan peran yang dipertukarkan), jadi pembatasannya normal:
        maka berinduksilah lalu rangkaikan basis eigen ortonormalnya. Sebaliknya,
        sebuah $u = UDU^\dagger$ yang \omterm{def:b2:reduction:diag}{dapat didiagonalkan} secara \omterm{def:b2:hermitian:adjoint}{uniter} memenuhi $u^*u
        = U\conj D D U^\dagger = U D\conj D U^\dagger = uu^*$:
        sehingga normal.
\end{enumerate}
\end{solution}

\begin{solution}{exo:b2:hermitian:9}
\omterm{def:b2:reduction:eigen}{Nilai eigennya}: untuk $u(x) = \lambda x$ dengan $x \neq 0$:
\[
\lambda\norm x^2 = \langle x, u(x)\rangle
= \langle u^*(x), x\rangle = -\langle u(x), x\rangle
= -\conj{\langle x, u(x)\rangle} = -\conj\lambda\,\norm x^2 ,
\]
jadi $\lambda = -\conj\lambda$: sehingga imajiner murni. Karena $(\iu
u)^* = -\iu\,u^*$ (sebab \omterm{def:b2:hermitian:adjoint}{adjoinnya} linear-sekawan pada skalarnya),
maka $u^* = u$ memberi $(\iu u)^* = -\iu u$: jadi pemetaan $u \mapsto \iu u$
mengirim yang \omterm{def:b2:hermitian:adjoint}{Hermitian} ke yang \omterm{def:b2:hermitian:adjoint}{anti-Hermitian}, dengan balikan $w \mapsto -\iu
w$: yakni sebuah bijeksi. Lalu $u$ yang \omterm{def:b2:hermitian:adjoint}{anti-Hermitian} memenuhi $u^*u = -u^2 =
uu^*$: jadi normal, sehingga \omterm{def:b2:reduction:diag}{dapat didiagonalkan} secara \omterm{def:b2:hermitian:adjoint}{uniter} menurut
\cref{exo:b2:hermitian:8}.
\end{solution}

\begin{solution}{exo:b2:hermitian:10}
\begin{enumerate}
  \item Pemetaan $S$ mempermutasikan sebuah basis ortonormal: sebab $\norm{Sx} = \norm
        x$, jadi $S$ \omterm{def:b2:hermitian:adjoint}{uniter}. Dengan mengindeks koordinatnya lewat $j = 0,
        \dots, n-1$ modulo $n$: $(Sx)_j = x_{j-1}$, jadi untuk
        $(f_k)_j = \frac{\omega^{jk}}{\sqrt n}$:
        \[
        (Sf_k)_j = \frac{\omega^{(j-1)k}}{\sqrt n}
        = \omega^{-k}\,(f_k)_j :
        \qquad Sf_k = \omega^{-k}f_k .
        \]
        Keortonormalannya: $\langle f_k, f_l\rangle = \frac1n
        \sum_j \omega^{j(l-k)} = \delta_{kl}$ (sebab jumlah geometri
        akar satuan yang tak trivial lenyap).
  \item $Cf_k = \sum_j c_j S^jf_k = \bigl(\sum_j
        c_j\omega^{-jk}\bigr)f_k = \widehat c(k)\,f_k$: jadi setiap
        matriks sirkulan bersifat diagonal dalam basis Fourier yang ortonormal,
        sehingga normal, dengan \omterm{def:b2:reduction:eigen}{spektrum} $\{\widehat c(k)\}$. (Dan
        penggantian basisnya adalah transformasi Fourier diskret:
        sebab konvolusi menjadi perkalian.)
\end{enumerate}
\end{solution}

\begin{solution}{exo:b2:hermitian:11}
($\Leftarrow$) Misalkan $P^2 = P = P^*$. Setiap $v$ terpecah menjadi $v = Pv
+ (v - Pv)$ dengan $Pv \in \operatorname{im} P$ dan $P(v - Pv) =
0$. Kedua kepingnya ortogonal: sebab untuk sembarang $x, y$,
\[
\langle Px, (I - P)y\rangle = \langle x, P(I-P)y\rangle
= \langle x, (P - P^2)y\rangle = 0 :
\]
jadi $\ker P \perp \operatorname{im} P$, sehingga $P$ adalah proyeksi
ortogonal ke petanya. ($\Rightarrow$) Bila $P$ adalah proyeksi
ortogonal ke $F = \operatorname{im}P$: maka untuk setiap
$x, y$ berlaku $\langle Px, y\rangle = \langle Px, Py\rangle$ (sebab
komponen $y - Py \perp F$ gugur) dan secara simetris $\langle
x, Py\rangle = \langle Px, Py\rangle$: jadi $\langle Px, y\rangle =
\langle x, Py\rangle$, yakni $P^* = P$. Matriksnya: ke $\C v$
(dengan $\norm v = 1$): $Px = v\,\langle v, x\rangle$, yakni $P =
vv^\dagger$; sedangkan ke $\operatorname{Vect}(v_1, \dots, v_k)$
yang ortonormal: $P = \sum_i v_iv_i^\dagger$.
\end{solution}

\begin{solution}{exo:b2:hermitian:12}
Diagonalkan $A = U D U^\dagger$ (lewat teorema spektral), dengan $D$ diagonal
berentri di antara $\lambda_i$. Maka $P_i = L_i(A) = U
L_i(D)U^\dagger$, dan $L_i(D)$ bersifat diagonal berentri
$L_i(\lambda_j) = \delta_{ij}$: yakni bernilai satu persis pada slot
$E_i$. Jadi $P_i$ bersifat \omterm{def:b2:hermitian:adjoint}{Hermitian} (sebab $L_i$ real dan $D$ real), idempoten,
berpeta $E_i$ dan berkernel $\bigoplus_{j\neq i}E_j = E_i^\perp$
(berkat keortogonalan \omterm{def:b2:reduction:eigen}{ruang eigennya}): sehingga ia proyeksi ortogonal ke
$E_i$ (\cref{exo:b2:hermitian:11}). Pola diagonal yang saling lepas
memberi $P_iP_j = 0$ (untuk $i \neq j$); lalu $\sum_i L_i = 1$ (sebab berderajat $< p$
dan bernilai $1$ di $p$ titik), jadi $\sum P_i = I$; dan $\sum_i
\lambda_iL_i(\lambda_j) = \lambda_j$ memberi $A = \sum
\lambda_iP_i$. Untuk sembarang polinomial $f$: matriks $f(D)$ berdiagonal
$f(\lambda_j)$, jadi
\[
f(A) = \sum_{i=1}^{p} f(\lambda_i)\,P_i :
\]
jadi fungsi $A$ dihitung secara spektral --- yakni kalkulus yang
diperluas jilid Tahun ke-3 ke $f$ yang \omterm{def:b2:metric:continuity}{kontinu} dan seterusnya.
\end{solution}

\begin{solution}{pb:b2:hermitian:1}
\textbf{1.} Berlaku $\conj{\langle x, Ax\rangle} = \langle Ax, x\rangle
= \langle x, A^*x\rangle = \langle x, Ax\rangle$: jadi real. Lalu dengan menulis
$x = \sum c_ie_i$:
\[
R_A(x) = \frac{\sum_i\lambda_i\abs{c_i}^2}
{\sum_i\abs{c_i}^2} ,
\]
yakni rata-rata terbobot \omterm{def:b2:reduction:eigen}{nilai eigennya}: jadi ia terletak di
$\intcc{\lambda_n}{\lambda_1}$, dengan batasnya tercapai di
$e_1$ dan $e_n$.

\textbf{2.} Untuk $t$ yang real dan sembarang $v$, uraikan $R_A(x + tv) =
\frac{N(t)}{D(t)}$ dengan
\[
N(t) = \langle x, Ax\rangle + 2t\Re\langle v, Ax\rangle +
t^2\langle v, Av\rangle,
\quad
D(t) = \norm x^2 + 2t\Re\langle v, x\rangle + t^2\norm v^2 .
\]
Turunannya di $t = 0$ adalah
\[
\frac{2}{\norm x^2}\,
\Re\bigl\langle v,\ Ax - R_A(x)\,x\bigr\rangle .
\]
Ia lenyap untuk setiap $v$ bila dan hanya bila $\Re\langle v, w\rangle = 0$ untuk
setiap $v$, dengan $w = Ax - R_A(x)x$; lalu mengganti $v$ dengan $\iu v$
membunuh bagian imajinernya juga: jadi $w = 0$, yakni $Ax = R_A(x)x$.
Sehingga titik kritis $R_A$ persis merupakan \omterm{def:b2:reduction:eigen}{vektor eigen}, dengan
nilai kritisnya berupa \omterm{def:b2:reduction:eigen}{nilai eigennya}.

\textbf{3.} Untuk $x = \sum_{i\leq k}c_ie_i \in V_k$: besaran $R_A(x)$
adalah rata-rata terbobot $\lambda_1, \dots, \lambda_k$, sehingga
$\geq \lambda_k$, dengan kesamaan di $e_k$: jadi $\min_{V_k} R_A =
\lambda_k$. Secara simetris pada $W_k$ rata-ratanya melibatkan
$\lambda_k, \dots, \lambda_n$: jadi $\max_{W_k}R_A = \lambda_k$.

\textbf{4.} Misalkan $\dim V = k$. Maka $\dim V + \dim W_k = n + 1 >
n$, jadi ada $x \in V \cap W_k$ yang satuan, dan $R_A(x) \leq
\lambda_k$ (pertanyaan 3): sehingga $\min_{V}R_A \leq \lambda_k$ untuk setiap
$V$ yang demikian. Karena $V_k$ mencapai $\lambda_k$, maka maks-minnya sama dengan
$\lambda_k$. Adapun rumus min-maksnya lewat hujah yang sama dengan
peran yang dibalik (sebab $\dim W = n - k + 1$ memaksa $W \cap V_k \neq
\{0\}$, jadi $\max_W R_A \geq \lambda_k$, yang tercapai di $W_k$).

\textbf{5.} Berlaku $R_B(x) = R_A(x) + \frac{\langle x,
(B-A)x\rangle}{\norm x^2} \geq R_A(x)$ \omterm{def:b2:funcseq:def}{titik demi titik}. Lalu mengambil
$\min$ atas sembarang ruang berdimensi $k$, sebut $V$, lalu $\max$ atas $V$ memberi
$\lambda_k(B) \geq \lambda_k(A)$ menurut pertanyaan 4.

\textbf{6.} Menurut Grassmann: $\dim(V\cap W) = \dim V + \dim W -
\dim(V + W) \geq \dim V + \dim W - n > 0$. Menerapkan ini dua kali:
\[
\dim(V_1\cap V_2\cap V_3)
\geq \dim(V_1\cap V_2) + \dim V_3 - n
\geq \dim V_1 + \dim V_2 + \dim V_3 - 2n .
\]

\textbf{7.} Subruang $W_i(A)$, $W_j(B)$ (pertanyaan 3, bagi
$A$ dan $B$) dan $V_{i+j-1}(A+B)$ berdimensi $(n-i+1) +
(n-j+1) + (i+j-1) = 2n + 1 > 2n$: jadi menurut pertanyaan 6 ada sebuah
vektor satuan $x$ di ketiganya. Maka
\[
\lambda_{i+j-1}(A+B) \leq R_{A+B}(x)
= R_A(x) + R_B(x) \leq \lambda_i(A) + \lambda_j(B),
\]
dengan ketaksamaan kirinya karena $x \in V_{i+j-1}(A+B)$ (pertanyaan
3), dan yang kanannya berkat kedua $W$-nya.

\textbf{8.} Dalam basis spektral $E$: $\norm{Ex}^2 = \sum
\lambda_k(E)^2\abs{c_k}^2 \leq \max_k\lambda_k(E)^2\,\norm
x^2$, yang tercapai di \omterm{def:b2:reduction:eigen}{vektor eigen} yang bersesuaian: jadi $\vertiii E_2
= \max_k\abs{\lambda_k(E)}$. Lalu Weyl dengan $j = 1$: $\lambda_k(A+E)
\leq \lambda_k(A) + \lambda_1(E) \leq \lambda_k(A) + \vertiii
E_2$; dan menerapkan ini pada $(A+E) + (-E)$: $\lambda_k(A) \leq
\lambda_k(A+E) + \vertiii E_2$. Bersama-sama:
$\abs{\lambda_k(A+E) - \lambda_k(A)} \leq \vertiii E_2$.

\textbf{9.} Batas bawahnya: berkat $P \geq 0$ dan pertanyaan 5. Batas atasnya: sebab $P$
berrank $1$, jadi $\lambda_2(P) = 0$; lalu Weyl dengan $i = k-1$ dan $j =
2$ memberi
\[
\lambda_k(A + P) \leq \lambda_{k-1}(A) + \lambda_2(P)
= \lambda_{k-1}(A) .
\]

\textbf{10.} Berlaku $A + E = \begin{pmatrix} 2 & 1+\iu\\ 1-\iu &
2\end{pmatrix}$: dengan \omterm{def:b2:reduction:charpoly}{polinomial karakteristik} $(2-\lambda)^2 -
\abs{1+\iu}^2 = (2-\lambda)^2 - 2$, jadi \omterm{def:b2:reduction:eigen}{spektrumnya} $\{2 + \sqrt2,\ 2
- \sqrt2\}$. Terhadap $\operatorname{Sp}A = \{3, 1\}$:
\[
\abs{(2+\sqrt2) - 3} = \abs{(2-\sqrt2) - 1} = \sqrt2 - 1
\approx 0.414 \leq 1 = \vertiii E_2 . \checkmark
\]

\textbf{11.} Pandanglah $\C^{n-1} = \operatorname{Vect}(e_1, \dots,
e_{n-1})$ di dalam $\C^n$ (lewat basis bakunya): untuk $x$ di sana,
$\langle x, Bx\rangle = \langle x, Ax\rangle$, jadi $R_B$ adalah
pembatasan $R_A$. Batas atasnya: maks-min bagi $\lambda_k
(B)$ berjangkau atas subruang berdimensi $k$ \emph{milik}
$\C^{n-1}$, yakni subkeluarga milik $\C^n$: jadi $\lambda_k(B) \leq
\lambda_k(A)$. Batas bawahnya: min-maks bagi $\lambda_k(B)$
berjangkau atas subruang $\C^{n-1}$ yang berdimensi $(n-1)-k+1 =
n-k$; dan masing-masingnya juga subruang $\C^n$ berdimensi $n -
(k+1) + 1$, jadi maksimumnya $\geq \lambda_{k+1}(A)$, sehingga
$\lambda_k(B) \geq \lambda_{k+1}(A)$.

\textbf{12.} Buanglah baris dan kolomnya satu demi satu lalu rantaikan
pertanyaan 11: sebab tiap pembuangannya menggeser indeks bawahnya satu, sehingga memberi
$\lambda_{k+m}(A) \leq \lambda_k(B) \leq \lambda_k(A)$.

\textbf{13.} Mengonjugasikan $A$ dengan matriks permutasi (yang \omterm{def:b2:hermitian:adjoint}{uniter})
tidak mengubah \omterm{def:b2:reduction:eigen}{spektrumnya} maupun multihimpunan entri
diagonalnya: jadi andaikan $d_1, \dots, d_k$ menempati posisi terdepannya.
Maka submatriks utama terdepan $k\times k$, sebut $B$, punya
$\operatorname{tr} B = \sum_{i\leq k}d_i$, dan \omterm{def:b2:reduction:eigen}{nilai eigennya}
memenuhi $\mu_i(B) \leq \lambda_i(A)$ (pertanyaan 12): jadi setelah dijumlahkan,
$\sum_{i\leq k}d_i \leq \sum_{i\leq k}\lambda_i(A)$. Dan di $k =
n$ kedua ruasnya adalah $\operatorname{tr} A$.

\textbf{14.} Mengambil $x_i = e_i$ memberi nilai
$\sum_{i\leq k}\lambda_i$: jadi maksimumnya $\geq$. Sebaliknya, sebuah
keluarga ortonormal $(x_1, \dots, x_k)$ diperluas menjadi
basis ortonormal, yakni menjadi $U$ yang \omterm{def:b2:hermitian:adjoint}{uniter} dengan kolom pertamanya
$x_i$; lalu $\sum_i\langle x_i, Ax_i\rangle$ adalah jumlah
$k$ entri diagonal pertama $U^\dagger AU$, yang menurut pertanyaan
13 paling banyak sebesar jumlah $k$ \omterm{def:b2:reduction:eigen}{nilai eigen} terbesarnya ---
yakni $\sum_{i\leq k}\lambda_i(A)$. Jadi asas maksimum Ky Fan
menyusul.

\textbf{15.} Penyisipannya (dengan $\operatorname{Sp}A = \{2+\sqrt2, 2,
2-\sqrt2\}$ dan $\operatorname{Sp}B = \{3, 1\}$):
\[
2 \leq 3 \leq 2 + \sqrt2,
\qquad
2 - \sqrt2 \leq 1 \leq 2 . \checkmark
\]
Schur dengan diagonal $(2,2,2)$: $2 \leq 2+\sqrt2$; lalu $4 \leq 4 +
\sqrt2$; dan $6 = 6$ (yakni tracenya). \checkmark

\textbf{16.} Berlaku $b_{ii} - a_{ii} = \langle e_i, (B-A)e_i\rangle
\geq 0$; dan menjumlahkannya memberi tracenya. Untuk sembarang $C$: $\langle x,
C^\dagger(B - A)Cx\rangle = \langle Cx, (B-A)(Cx)\rangle \geq
0$: jadi $C^\dagger AC \leq C^\dagger BC$.

\textbf{17.} Jelas $A \geq 0$; dan $B - A = \begin{pmatrix} 1 & 1\\
1 & 1\end{pmatrix}$ bersifat semidefinit positif (bernilai eigen $2,
0$): jadi $A \leq B$. Tetapi
\[
B^2 = \begin{pmatrix} 5 & 3\\ 3 & 2\end{pmatrix},
\qquad
B^2 - A^2 = \begin{pmatrix} 4 & 3\\ 3 & 2\end{pmatrix},
\qquad \det(B^2 - A^2) = -1 < 0 :
\]
sehingga tak semidefinit positif. Jadi pengkuadratan tak menghormati
urutan Loewner.

\textbf{18.} Misalkan $S = \sqrt A$ dan $T = \sqrt B$ (yang \omterm{def:b2:hermitian:adjoint}{Hermitian}
semidefinit positif, sebab \cref{exo:b2:hermitian:7} diperluas ke yang
semidefinit lewat rumus spektral yang sama). Maka $T - S$ bersifat
\omterm{def:b2:hermitian:adjoint}{Hermitian}; misalkan $\mu$ sembarang \omterm{def:b2:reduction:eigen}{nilai eigennya} dan $v$ \omterm{def:b2:reduction:eigen}{vektor eigen} satuannya.
Dari $T^2 - S^2 = T(T - S) + (T - S)S$:
\[
0 \leq \langle v, (B - A)v\rangle
= \langle Tv, (T-S)v\rangle + \langle (T-S)v, Sv\rangle
= \mu\bigl(\langle v, Tv\rangle + \langle v,
Sv\rangle\bigr)
\]
(sebab $\mu$ real). Bila $\langle v, Tv\rangle + \langle v,
Sv\rangle > 0$, maka $\mu \geq 0$. Bila ia lenyap, kedua suku
taknegatifnya lenyap; lalu $\langle v, Tv\rangle =
\norm{T^{1/2}v}^2 = 0$ memaksa $Tv = 0$, dan demikian pula $Sv = 0$, jadi
$\mu v = (T - S)v = 0$ sehingga $\mu = 0$. Jadi semua \omterm{def:b2:reduction:eigen}{nilai eigen} $T -
S$ bernilai $\geq 0$: sehingga $\sqrt A \leq \sqrt B$.

\textbf{19.} \omterm{def:b2:quadratic:def}{Kekongruenan} oleh $A^{-1/2}$ (pertanyaan 16) memberi $I \leq M
:= A^{-1/2}BA^{-1/2}$. Jadi semua \omterm{def:b2:reduction:eigen}{nilai eigen} $M$ bernilai $\geq 1$,
sehingga \omterm{def:b2:reduction:eigen}{nilai eigen} $M^{-1}$ terletak di $\intoc{0}{1}$: jadi $M^{-1} \leq I$.
Tetapi $M^{-1} = A^{1/2}B^{-1}A^{1/2}$; jadi mengongruenkan $M^{-1} \leq
I$ dengan $A^{-1/2}$ memberi $B^{-1} \leq A^{-1}$.

\textbf{20.} Berlaku $\sqrt A^{-1}(AB)\sqrt A = \sqrt A\,B\sqrt A$, jadi
$AB$ serupa dengan matriks \omterm{def:b2:hermitian:adjoint}{Hermitian} definit positif $\sqrt
A\,B\sqrt A$ (yang definit sebab $\langle x, \sqrt AB\sqrt Ax\rangle =
\langle \sqrt Ax, B\sqrt Ax\rangle > 0$): sehingga \omterm{def:b2:reduction:eigen}{nilai eigennya}
real dan positif. Lebih jauh
\[
\langle x, \sqrt AB\sqrt Ax\rangle
\leq \lambda_{\max}(B)\,\norm{\sqrt Ax}^2
= \lambda_{\max}(B)\,\langle x, Ax\rangle
\leq \lambda_{\max}(A)\lambda_{\max}(B)\norm x^2 ,
\]
jadi $\lambda_{\max}(AB) = \max R_{\sqrt AB\sqrt A} \leq
\lambda_{\max}(A)\lambda_{\max}(B)$.

\textbf{21.} Dengan $v_k = (\sin j\theta_k)_j$ dan kesamaan
$\sin((j-1)\theta) + \sin((j+1)\theta) =
2\sin(j\theta)\cos\theta$: untuk $2 \leq j \leq n-1$,
\[
(T_nv_k)_j = -\sin((j{-}1)\theta_k) + 2\sin(j\theta_k) -
\sin((j{+}1)\theta_k)
= (2 - 2\cos\theta_k)\sin(j\theta_k) .
\]
Baris $1$ berhasil karena $\sin(0\cdot\theta_k) = 0$, dan baris $n$
karena $\sin((n+1)\theta_k) = \sin(k\pi) = 0$: jadi syarat
perbatasannya memilih persis $\theta_k = \frac{k\pi}{n+1}$. Sehingga
$T_nv_k = 4\sin^2\bigl(\frac{k\pi}{2(n+1)}\bigr)v_k$; adapun $n$
nilainya berbeda-beda di $\intoo04$ dan $v_k \neq 0$: jadi inilah
seluruh \omterm{def:b2:reduction:eigen}{spektrumnya}, yang positif, sehingga $T_n$ definit positif.

\textbf{22.} Berlaku $\min_id_i\,I \leq D \leq \max_id_i\,I$, jadi $T_n +
\min_id_i\,I \leq T_n + D \leq T_n + \max_id_i\,I$
(sebab menambahkan $T_n$ memelihara urutannya); lalu pertanyaan 5 dan
$\lambda_k(T_n + cI) = \lambda_k(T_n) + c$ memberi apitannya.

\textbf{23.} Matriks $T_2 = \begin{pmatrix} 2 & -1\\ -1 &
2\end{pmatrix}$ berspektrum $\{3, 1\}$, dan
\[
2 - \sqrt2 \;\leq\; 1 \;\leq\; 2 \;\leq\; 3 \;\leq\; 2 +
\sqrt2 :
\]
jadi penyisipan Cauchynya terperiksa. (Dan inilah \omterm{def:b2:reduction:eigen}{spektrum} yang sama seperti pada
pertanyaan 15: sebab mengonjugasikan dengan $\operatorname{diag}(1,-1,1)$ membalik
tanda luar diagonalnya.) Pembacaan grafnya: $T_2$ adalah
matriks bertipe Laplacian bagi lintasan yang simpul terakhirnya dihapus
--- yakni sebuah submatriks utama, persis situasi pertanyaan
11.

\textbf{24.} \omterm{def:b2:reduction:eigen}{Nilai eigen} ekstremnya berkelakuan sebagai
\[
\lambda_{\min}(T_n) = 4\sin^2\frac{\pi}{2(n+1)}
\sim \frac{\pi^2}{(n+1)^2},
\qquad
\lambda_{\max}(T_n) = 4\cos^2\frac{\pi}{2(n+1)}
\longrightarrow 4 .
\]
Karena itu
\[
\kappa_n = \frac{\lambda_{\max}}{\lambda_{\min}}
\sim \frac{4(n+1)^2}{\pi^2} :
\]
jadi makin halus kisinya, makin buruk kondisi turunan kedua
diskretnya --- yakni fakta yang menggerakkan perancangan aljabar linear
numerik.

\textbf{25.} (i) Akar \omterm{def:b2:reduction:charpoly}{polinomial karakteristik} dapat
bergerak liar di bawah usikan sebuah matriks umum, tetapi
pencirian min-maksnya memaku tiap \omterm{def:b2:reduction:eigen}{nilai eigen} \omterm{def:b2:hermitian:adjoint}{Hermitian} antara
nilai pengoptimalan yang gamblang, sehingga memaksa kestabilan
\omterm{def:b2:metric:continuity}{Lipschitz} berkonstanta $1$ pada pertanyaan 8--9. (ii) Pertanyaan 1, 5, 16--20 hanya
memakai nilai Rayleigh ekstremnya; sedangkan Weyl, penyisipan, Schur dan
Ky Fan (pertanyaan 7--14) sungguh menuntut min-maks selengkapnya
atas subruangnya. (iii) Urutan Loewner mengubah ketaksamaan skalar
itu menjadi kalkulus ketaksamaan matriks --- dengan
jebakan yang nyata (pertanyaan 17) dan teorema yang nyata (pertanyaan 18--19).
(iv) Perkakas ini adalah roti sehari-hari analisis numerik
(yakni matriks kekakuan pertanyaan 24), teori usikan
kuantum (yakni Weyl: aras energinya bergeser paling banyak sebesar \omterm{def:b2:nvs:norm}{norma}
usikannya), dan asas min-maks pada jilid Tahun ke-3
bagi operator \omterm{def:b2:hermitian:adjoint}{swa-adjoin} yang \omterm{def:b2:metric:compact}{kompak}.

\end{solution}
